\section{Immutable Parent Pointers}
\label{sec:immutable-parent-pointers}

The immutable parent pointer is the foundational structural primitive of the Recursive Spawning Protocol. It is the single mechanism that transforms the delegation chain from a forensic reconstruction — assembled retroactively from mutable event logs and policy assertions — into a runtime-verifiable artifact: immutable, cryptographically enforced, and architecturally verifiable at any execution point. Every RSP guarantee — drift prevention, chain integrity, termination, and scope refinement — derives structurally from this one primitive. This section specifies the parent pointer formally, defines the chain traversal operator, establishes the root human anchor as the chain base case, describes immutability enforcement through the Fact Layer write protocol, explains the Purpose Layer's exclusive role in spawn authorization, and provides the complete chain traversal algorithms with correctness arguments.

% ------------------------------------------------------------
\subsection{Formal Definition: The ParentPointer Relation}
\label{sec:parent-pointer-formal}

\begin{dfn}[Parent Pointer]\label{dfn:parent-pointer}
Let $\mathcal{N}$ be the set of all agent nodes in a \bxthree{} system implementing the Recursive Spawning Protocol, and let $\mathcal{H} \subset \mathcal{N}$ denote the subset of nodes designated as human principal anchors. For any child node $c \in \mathcal{N}$ that was spawned via a valid RSP spawn event, the \emph{parent pointer} $\ParentPointer{p}{c}$ is the immutable, cryptographically signed reference from child $c$ to its immediate parent $p \in \mathcal{N}$, established exactly once at node provisioning via a Fact Layer atomic write operation.

The parent pointer satisfies four structural constraints for all nodes in $\mathcal{N}$:

\begin{enumerate}
    \item[(P1) \textbf{Uniqueness}] For every spawned node $c \in \mathcal{N}$, there exists exactly one parent $p \in \mathcal{N}$ such that $\ParentPointer{p}{c}$ holds after provisioning. No node may have multiple simultaneous parent pointers. No parent pointer may be removed, redirected, or overwritten after establishment. The relation is total and single-valued on $\mathcal{N} \setminus \mathcal{H}$.

    \item[(P2) \textbf{Immutability}] After the provisioning write confirming $\ParentPointer{p}{c}$ is recorded in the Fact Layer forensic ledger, no operation exists — from any source, with any privilege level, under any system condition including administrative compromise, software vulnerability, or hardware failure — that can modify or delete $\ParentPointer{p}{c}$. The pointer is permanent for the operational lifetime of $c$. This is not a policy enforced by access control; it is a structural property of the protocol. The modification operation is not defined in the Fact Layer instruction set.

    \item[(P3) \textbf{Directionality}] The parent pointer is an intrinsic attribute of the child node, recorded in the child's own Fact Layer state record. The parent node $p$ holds no mutable reference to $c$; the parent does not ``own'' its children through a held pointer. This means the chain is traversable from any node outward to its complete lineage regardless of the current operational state of any other node in the chain. A parent that has terminated, crashed, or been decommissioned does not affect the child's ability to reference its parent through $\ParentPointer{\cdot}{c}$.

    \item[(P4) \textbf{Human Anchor Termination}] The root of any RSP chain is a human principal node $h \in \mathcal{H}$ for which $\ParentPointer{\bot}{h}$ holds — the null parent marker indicating that $h$ is the ultimate delegation origin and the chain terminates at a named human. No machine actor in $\mathcal{N} \setminus \mathcal{H}$ may serve as a chain root.
\end{enumerate}
\end{dfn}

\medskip
\noindent\textbf{Notation.} We write $\ParentPointer{p}{c}$ to denote the parent pointer relation between parent $p$ and child $c$. We write $\alpha(c) = p$ as equivalent functional notation. The expression $\alpha(c) = \bot$ denotes that node $c$ has no parent — it is the root human anchor. We write $\parent(c)$ to denote the unique $p$ such that $\ParentPointer{p}{c}$ holds, and $c \in \children(p)$ for the inverse relation.

\medskip
\noindent\textbf{Why immutability is not an access-control problem.} The distinction between access-control-based pointer immutability and the Fact Layer write protocol is the central architectural difference that makes RSP guarantees structural rather than conventional. In a conventional access-control model, immutability is enforced by restricting which principals may issue modification operations. An actor who acquires sufficient privileges — through credential theft, privilege escalation, or administrative override — can modify the pointer. The immutability was a policy, enforceable until circumvention became possible.

In the Fact Layer write protocol, immutability is a structural property: the modification operation does not exist in the instruction set. No amount of privilege elevation, credential compromise, or software vulnerability can produce a valid modification operation against $\ParentPointer{p}{c}$ after provisioning. The protocol does not enforce immutability against adversarial actors — it renders post-provisioning modification architecturally incoherent. This is not a stronger access control policy; it is a different category of guarantee.

% ------------------------------------------------------------
\subsection{The Chain Operator: Recursive Definition and Base Case}
\label{sec:chain-operator}

The chain operator $\Chain{\cdot}$ is the primary mechanism for chain traversal, accountability attribution, and runtime verification. It constructs the complete delegation lineage from any node to the human anchor by recursively applying the parent pointer relation.

\begin{dfn}[Chain Operator]\label{dfn:chain-operator}
For any agent node $a \in \mathcal{N}$ in a Recursive Spawning chain, the \emph{chain} $\Chain{a}$ is defined recursively as:

\begin{equation}
\Chain{a} =
\begin{cases}
\{\, a \,\} \cup \Chain{\parent(a)} & \text{if } \alpha(a) \neq \bot \\[12pt]
\{\, a \,\} & \text{if } \alpha(a) = \bot \quad \text{(root human anchor)}
\end{cases}
\label{eq:chain-recursive}
\end{equation}

Equivalently, expanding the recursion:

\begin{equation}
\Chain{a} = \bigl\{\, a,\ \alpha(a),\ \alpha^{(2)}(a),\ \alpha^{(3)}(a),\ \ldots,\ \alpha^{(k)}(a) \,\bigr\}
\label{eq:chain-expanded}
\end{equation}

where $k$ is the chain depth of $a$, $\alpha^{(j)}(a)$ denotes the $j$-fold application of $\alpha$ (with $\alpha^{(0)}(a) = a$), and $\alpha^{(k)}(a)$ is the unique root node $n_0$ satisfying $\alpha(n_0) = \bot$ and $n_0 \in \mathcal{H}$.
\end{dfn}

\begin{cor}[Chain Termination]\label{cor:chain-termination}
Every accountability chain $\Chain{a}$ terminates at a node $h$ with $\alpha(h) = \bot$. By constraint (P4) of Definition~\ref{dfn:parent-pointer}, $h \in \mathcal{H}$; the terminus of every chain is a human principal. No RSP chain terminates at a machine actor.
\end{cor}

\medskip
\noindent\textbf{Base case: root human anchor.} The root of every RSP chain is a Purpose Layer entity designated as a human principal $h \in \mathcal{H}$, established at system initialization when $h$ authorizes the first agent node $n_0$. The Purpose Layer records $h(n_0) = h$ and $\alpha(n_0) = \bot$ in the Fact Layer authorization ledger. This is the only valid chain-root in RSP: no machine actor may serve as a chain origin, and no node may have $\alpha(n) = \bot$ unless it carries the $\hmannchor{\cdot}$ designation. The non-collapsing human anchor property ensures $h(n) = h(n_0)$ for all nodes in all chains spawned from $n_0$.

The base case terminates the recursion. Because every RSP chain has finite depth bounded by $D_{\max}$, the recursion in Equation~\ref{eq:chain-recursive} terminates at the root within at most $D_{\max} + 1$ successive applications of $\alpha$.

\medskip
\noindent\textbf{Correctness of the recursive definition.} The recursive definition in Equation~\ref{eq:chain-recursive} is correct by structural induction on chain depth. The base case holds for depth-0 nodes (human anchors). For the inductive case, assume $\Chain{\parent(a)}$ is correctly defined for all nodes with depth less than $\text{depth}(a)$. Since $\alpha(a) \neq \bot$, $a$ has a parent with strictly smaller depth. By the inductive hypothesis, $\Chain{\parent(a)}$ is the correct chain from the parent to the human anchor. Prepending $a$ to this set yields precisely the chain from $a$ to the human anchor. By induction, Equation~\ref{eq:chain-recursive} holds for all nodes.

\begin{invariant}[Chain Integrity]\label{inv:chain-integrity}
For any node $c \in \mathcal{N}$ with $\ParentPointer{p}{c}$, we have $p \in \Chain{c}$ and $p$ appears immediately after $c$ in the ordered chain. Equivalently: $\Chain{c} = \{c\} \cup \Chain{p}$.
\end{invariant}

This invariant follows directly from Definition~\ref{dfn:chain-operator}. It guarantees the chain is a proper sequence: no node appears twice, and no links are skipped in the delegation path.

% ------------------------------------------------------------
\subsection{Immutability Enforcement: The Fact Layer Write Protocol}
\label{sec:fact-layer-write-protocol}

The Fact Layer enforces parent pointer immutability through a write protocol that defines no valid modification operation for $\ParentPointer{p}{c}$ after provisioning. This is not access control with elevated privileges — it is the absence of a modification pathway.

The Fact Layer write protocol enforces four structural properties that collectively make immutability a property of the protocol itself:

\begin{enumerate}
    \item[(I1) \textbf{No defined modification operation}] The Fact Layer instruction set does not include an operation to overwrite, redirect, or delete $\ParentPointer{p}{c}$ after the provisioning write is confirmed. The operation does not exist in the protocol specification. This is not a restriction on who may call a modification operation — it is a statement that the operation is not defined.

    \item[(I2) \textbf{No override path for any actor}] There is no administrative override, no emergency bypass, no operational urgency flag, and no privileged actor — including the Purpose Layer after provisioning, the Bounds Engine, system security officers, or any external principal — that can issue a valid parent pointer modification. Modification requests from any source receive structurally identical rejection. There is no privileged source.

    \item[(I3) \textbf{Failure-state resilience}] The immutability guarantee holds under all system failure conditions, including partial failures that disable other enforcement mechanisms. The Fact Layer maintains its own state integrity independently of the operational state of other system components. A node whose parent has terminated, crashed, or been decommissioned still has an immutable $\alpha$ value pointing to that parent.

    \item[(I4) \textbf{Atomic warrant-pointer provisioning}] The parent pointer and its Purpose Layer warrant are created as a single atomic Fact Layer write. There is no temporal window in which $\ParentPointer{p}{c}$ exists without a matching warrant, or in which a warrant exists without a corresponding parent pointer. This eliminates the class of exploits in which a parent pointer is first established and then retroactively authorized after the fact.
\end{enumerate}

\medskip
\noindent\textbf{The sealed memory envelope.} The immutability guarantee is physically reinforced by the sealed memory envelope mechanism:

\begin{enumerate}
    \item \textbf{Encrypted at rest.} The parent pointer field $\alpha(c)$ is encrypted in the node's persistent memory envelope using a Fact Layer-only symmetric key. Even a compromised Bounds Engine cannot read $\alpha(c)$ because it never has access to the decryption key.

    \item \textbf{HMAC-signed.} The envelope carries an HMAC-SHA-256 signature computed over its contents using a Fact Layer-only key. Any modification to the ciphertext invalidates the signature.

    \item \textbf{Decrypted only for reads.} The Fact Layer decrypts the envelope only when reading $\alpha(c)$ to service a chain-traversal or audit request. The plaintext is never exposed to the Bounds Engine or the child node's own code.

    \item \textbf{Re-sealed after reads.} After each authorized read, the Fact Layer re-encrypts and re-signs the envelope, preventing replay attacks based on stale plaintext.
\end{enumerate}

The Fact Layer operates as a standalone, isolated process with its own memory space, inaccessible to the child node's address space. The combination of the protocol-level write constraint and the sealed envelope ensures that $\alpha(c)$ satisfies the modification-resistance property across all deployment contexts.

\medskip
\noindent\textbf{The Fact Layer pre-commit hook.} Every spawn event passes through the Fact Layer pre-commit hook before the provisioning write is recorded. The hook verifies two structural invariants for each spawn event $n_i \to n_{i+1}$:

\begin{enumerate}
    \item[(FC1)] $\ParentPointer{n_i}{n_{i+1}}$ is being established with a Purpose Layer warrant present in the ledger — confirming that a machine actor cannot self-originate a delegation chain.

    \item[(FC2)] The depth constraint $\text{depth}(n_{i+1}) \leq D_{\max}$ is re-verified at the Fact Layer even though it was already checked at the Bounds Engine — providing a redundant enforcement point that prevents drift through single-point failures.
\end{enumerate}

If either invariant fails, the hook rejects the spawn and logs the rejection. The pre-commit hook is the structural mechanism that prevents drift conditions from being recorded — it enforces constraints at the write level, not the policy level.

\begin{prop}[Immutability of ParentPointer]\label{prop:parent-pointer-immutable}
For any child node $c$ with parent $p$, once the spawn event $\mathrm{spawn}(c, p, \sigma_c, t, \phi)$ has been recorded in the forensic ledger $\mathcal{L}$, the value $\alpha(c)$ is fixed for the operational lifetime of $c$. No subsequent operation — including operations issued by $p$, by any Purpose Layer actor, by the Bounds Engine, or by $c$ itself — can alter $\alpha(c)$.
\end{prop}

This property follows directly from (I1)–(I4) and distinguishes ParentPointer from conventional mutable parent references in multi-agent systems. A mutable reference can be overwritten by a subsequent administrative operation; an immutable pointer cannot. The distinction is not behavioral — it is architectural.

% ------------------------------------------------------------
\subsection{Spawn Authorization: Purpose Layer Role}
\label{sec:purpose-layer-authorization}

The Purpose Layer plays two structurally necessary roles in the parent pointer lifecycle: it defines the spawn authorization policy that governs when a new parent pointer may be created, and it issues the spawn warrant that is the prerequisite for the Fact Layer provisioning write. These roles are exclusive to the Purpose Layer by architectural necessity — if a machine actor could create a parent pointer without a Purpose Layer warrant, or if the Purpose Layer could issue warrants retroactively, the chain's accountability guarantees would collapse.

\medskip
\noindent\textbf{Authorization policy origination.} At system initialization, the Purpose Layer defines the spawn authorization policy $P_{\text{spawn}}$ and records it in the Fact Layer authorization ledger. $P_{\text{spawn}}$ specifies the mandatory conditions for any valid spawn event:

\begin{enumerate}
    \item[(S1)] \textbf{Scope refinement.} The child scope must be a strict subset of the parent scope: $R(n_{\text{child}}) \subset R(n_{\text{parent}})$. No lateral expansion, no superset, no equality.

    \item[(S2)] \textbf{Operational necessity.} The parent must declare, in the Fact Layer authorization ledger, the precise condition requiring child spawning. The declaration must identify why the condition cannot be resolved within $R(n_{\text{parent}})$ and why the child is the minimum necessary decomposition of the current task.

    \item[(S3)] \textbf{Minimum scope proportionality.} The child scope must be the narrowest scope sufficient to address the declared condition. The Purpose Layer does not authorize child scopes broader than operationally necessary.

    \item[(S4)] \textbf{Human root confirmation.} The parent must demonstrate that the complete chain from the proposed child to the human anchor $h \in \mathcal{H}$ is intact and verified — that every node in the chain is responsive and that the chain terminates at a human principal.
\end{enumerate}

\medskip
\noindent\textbf{Warrant issuance.} When all conditions in $P_{\text{spawn}}$ are satisfied, the Purpose Layer issues a spawn warrant $w_i$ — a cryptographically signed record written atomically to the Fact Layer authorization ledger:

\begin{equation}
w_i = \bigl\langle
\text{parent}=n_i,\;
\text{child}=n_{i+1},;
R(n_{i+1}),;
\text{justification},;
\text{timestamp},;
\sigma_{\text{PL}}
\bigr\rangle .
\label{eq:spawn-warrant}
\end{equation}

The warrant is the sole authorized precondition for spawn. No machine actor may initiate a spawn event without a matching warrant in the Fact Layer ledger — the Fact Layer pre-commit hook rejects any provisioning request that lacks a valid warrant before the operation reaches the ledger write stage. This makes the Purpose Layer's exclusive terminus role a structural guarantee, not a configuration option.

\medskip
\noindent\textbf{Human anchor establishment.} At root provisioning — when a human principal $h$ first authorizes the agent node $n_0$ — the Purpose Layer records the human anchor $h(n_0) = h$ in the Fact Layer authorization ledger. The anchor is non-collapsing: for all nodes $n_j$ in any chain descended from $n_0$, $h(n_j) = h(n_0) = h$. The anchor does not change as the chain deepens. No machine actor can serve as a chain root.

The separation of authorization (Purpose Layer) from assignment (Fact Layer) is essential: it prevents a parent from claiming to have spawned a node whose parent pointer was set to a different actor, and it prevents a child from falsifying its own ancestry after provisioning.

% ------------------------------------------------------------
\subsection{Chain Traversal Algorithm}
\label{sec:chain-traversal-algorithm}

The chain traversal algorithm is the runtime procedure for constructing the delegation chain from any node to its human anchor. It is used by the Bailout Protocol for exception escalation, by auditors for accountability verification, and by the Fact Layer's continuous orphan detection. It provides the operational counterpart to the chain operator definition.

\medskip
\noindent\textbf{Algorithm: Chain-Traverse}

\begin{algorithm}
\caption{Chain-Traverse($n$) — Build the delegation chain from node $n$ to the human anchor}
\label{alg:chain-traverse}
\begin{algorithmic}[1]
\REQUIRE $n$ is a provisioned RSP node with $\alpha$ defined
\ENSURE $\Chain{n}$ — the ordered set of all ancestor nodes including the human anchor

\STATE $chain \leftarrow \langle\ \rangle$ \COMMENT{Initialize empty ordered list}
\STATE $current \leftarrow n$ \COMMENT{Start at the query node}
\WHILE{$\text{true}$}
    \STATE $\text{append}(chain,\ current)$ \COMMENT{Add current node to chain}
    \IF{$\alpha(current) = \bot$}
        \STATE $\text{append}(chain,\ current)$ \COMMENT{Add human anchor and terminate}
        \STATE \textbf{break} \COMMENT{Root human anchor reached}
    \ENDIF
    \STATE $current \leftarrow \alpha(current)$ \COMMENT{Move to parent}
\ENDWHILE

\RETURN $chain$
\end{algorithmic}
\end{algorithm}

\medskip
\noindent\textbf{Correctness arguments for Chain-Traverse.}

\textit{Termination.} The algorithm terminates because RSP chains are depth-bounded by construction (proved in Section~\ref{sec:rs-correctness}). Each iteration of the while-loop moves $\text{current}$ to the parent node $\alpha(current)$, which is strictly closer to the root along the chain. Since the chain has finite depth at most $D_{\max} + 1$, the loop executes at most $D_{\max} + 1$ iterations before $\alpha(current) = \bot$ is reached. This argument holds regardless of the current operational state of any node in the chain, including terminated or crashed nodes, because the traversal only reads immutable $\alpha$ values.

\textit{Correctness of the returned chain.} By Definition~\ref{dfn:chain-operator}, $\Chain{a}$ is the set containing $a$, $\alpha(a)$, $\alpha(\alpha(a))$, and so on, until the root. The algorithm produces exactly this sequence: it starts at $a$, appends each node, moves to its parent via $\alpha$, and stops when $\alpha(current) = \bot$ is reached. By structural induction on the number of loop iterations, the output is identical to $\Chain{a}$.

\textit{Uniqueness of result.} The algorithm produces a unique result for any input node because $\alpha$ is a total function (defined for every provisioned node) and is injective on the chain direction (every node has exactly one parent; the root has no parent). There is no branching, no alternative path, and no non-deterministic choice at any step.

\textit{Immutability guarantee.} The algorithm reads $\alpha$ values but never writes them. Traversal is read-only and does not interact with the Fact Layer write protocol. The algorithm can be executed at any runtime point without affecting the immutability of the chain it traverses.

\medskip
\noindent\textbf{Algorithm: Chain-Verify}

Chain-Verify confirms the structural validity of a delegation chain at any runtime point. It is the runtime verification procedure corresponding to the chain operator definition — confirming that a given node's chain satisfies all RSP structural constraints.

\begin{algorithm}
\caption{Chain-Verify($n$) — Verify that node $n$ terminates at a human anchor via a valid RSP chain}
\label{alg:chain-verify}
\begin{algorithmic}[1]
\REQUIRE $n$ is a provisioned RSP node
\ENSURE $\text{Boolean}$ indicating whether the chain is a structurally valid RSP delegation chain

\STATE $chain \leftarrow \text{Chain-Traverse}(n)$
\STATE $m \leftarrow \text{len}(chain)$

\STATE \COMMENT{\textit{Step V1: Verify chain terminates at human anchor}}
\IF{$\alpha(chain[m-1]) \neq \bot$}
    \STATE \RETURN $\text{false}$ \COMMENT{Root does not have null parent — malformed chain}
\ENDIF
\IF{$\neg \hmannchor{chain[m-1]}$}
    \STATE \RETURN $\text{false}$ \COMMENT{Root is not a designated human principal}
\ENDIF

\STATE \COMMENT{\textit{Step V2: Verify Purpose Layer warrant exists for each spawn}}
\FOR{$i \leftarrow 0$ \TO $m-2$}
    \STATE $child \leftarrow chain[i]$
    \STATE $parent \leftarrow chain[i+1]$
    \IF{$\neg \text{warrant\_exists}(child,\ parent)$}
        \STATE \RETURN $\text{false}$ \COMMENT{Missing Purpose Layer warrant}
    \ENDIF
\ENDFOR

\STATE \COMMENT{\textit{Step V3: Verify scope refinement at each link}}
\FOR{$i \leftarrow 0$ \TO $m-2$}
    \STATE $child \leftarrow chain[i]$
    \STATE $parent \leftarrow chain[i+1]$
    \IF{$\neg (R(child) \subset R(parent))$}
        \STATE \RETURN $\text{false}$ \COMMENT{Scope refinement violated}
    \ENDIF
\ENDFOR

\STATE \COMMENT{\textit{Step V4: Verify depth bound compliance}}
\FOR{$i \leftarrow 0$ \TO $m-1$}
    \IF{$\text{depth}(chain[i]) > D_{\max}$}
        \STATE \RETURN $\text{false}$ \COMMENT{Depth bound exceeded}
    \ENDIF
\ENDFOR

\RETURN $\text{true}$
\end{algorithmic}
\end{algorithm}

Chain-Verify returns \texttt{true} if and only if all four verification conditions hold simultaneously. V1 confirms the chain terminates at a human anchor — not at a machine actor or orphaned node. V2 confirms every spawn was authorized by the Purpose Layer — not self-bootstrapped by a machine actor. V3 confirms scope refinement was respected at every link — not lateral expansion. V4 confirms the depth bound was respected throughout. A chain that passes Chain-Verify is, by RSP definitions, a fully authorized delegation chain terminating at a human principal.

The algorithm is deterministic and produces the same result regardless of which node in the chain initiates verification. There is no self-serving interpretation of authorization that a drifted or orphaned node could produce, because Chain-Verify inspects the complete chain topology and warrants in the Fact Layer ledger — not the node's own claims about its authorization.

\medskip
\noindent\textbf{Computational cost.} Chain-Traverse runs in $O(d)$ time where $d$ is chain depth, bounded by $D_{\max} + 1$. Chain-Verify runs in $O(m)$ time where $m$ is chain length, plus three additional linear passes. Both algorithms are linear in chain length and do not depend on the total number of nodes in the system. The worst-case number of pointer dereferences to reach the human anchor from any node is $D_{\max} + 1$, independent of system scale.

The constant-time escalation bound is critical for the Bailout Protocol, where worst-case exception escalation latency must be bounded and predictable. A system with 10,000 active nodes and $D_{\max} = 5$ requires at most six pointer dereferences to reach the human anchor — identical to a system with 10 nodes. Mutable parent pointers would allow the chain to be extended retroactively, making the effective depth potentially unbounded and the worst-case escalation time unbounded as well. Immutability is what makes the depth bound a structural guarantee, not merely a policy parameter.

In the Bailout Protocol context, chain traversal identifies the correct Human Accountability Anchor when an agent at any depth encounters an unresolvable exception. The upward propagation of the exception signal follows the same parent-pointer structure as chain traversal, but uses the asynchronous trigger pathway rather than the synchronous verification pathway. Chain traversal ensures this walk never diverges, never shortcuts, and never fails to reach a human.
